\documentclass[reqno, 12pt]{amsart}
\usepackage[brazil]{babel}
\usepackage[utf8]{inputenc}
\usepackage[T1]{fontenc}
\usepackage{lmodern}
\usepackage{amsfonts,amssymb,amsmath,textcomp,amsthm}
\usepackage{tikz}
\usetikzlibrary{calc}
\usepackage{geometry}
\geometry{left=25mm, right=25mm, top=22mm, bottom=28mm}
\usepackage{setspace}
\onehalfspacing
\usepackage{csquotes}
\usepackage[hidelinks]{hyperref}
\usepackage[backend=biber, style=numeric, maxbibnames=99]{biblatex}
\addbibresource{refs_antiramsey.bib}
\setlength{\parindent}{1.2cm}
\numberwithin{equation}{section}

% -------------------------------------------------------
% Ambientes
% -------------------------------------------------------
\theoremstyle{plain}
\newtheorem{theorem}{Teorema}[section]
\newtheorem{proposition}[theorem]{Proposição}
\newtheorem{lemma}[theorem]{Lema}
\newtheorem{corollary}[theorem]{Corolário}
\newtheorem{afirmacao}[theorem]{Afirmação}

\theoremstyle{definition}
\newtheorem{definition}[theorem]{Definição}
\newtheorem{question}[theorem]{Questão}
\newtheorem{problem}[theorem]{Problema}
\newtheorem{remark}[theorem]{Observação}
\newtheorem{example}[theorem]{Exemplo}

% -------------------------------------------------------
% Macros
% -------------------------------------------------------
\newcommand{\Gnp}{\mathcal{G}(n,p)}
\newcommand{\rb}{\xrightarrow{\,\mathrm{rb}\,}}
\newcommand{\ramsey}[1]{\xrightarrow{\,#1\,}}
\newcommand{\cram}{\xrightarrow{\,\mathrm{c\text{-}ram}\,}}
\newcommand{\mdens}{m_2}
\newcommand{\maxdens}{m}
\newcommand{\whp}{a.a.s.}
\newcommand{\hugo}[1]{\textcolor{blue}{H: #1}}

\date{\today}

% ============================================================
\begin{document}
% ============================================================

\begin{center}
  {\large\textbf{Ramsey Assimétrico e Anti-Ramsey em Grafos Aleatórios}}

  \vspace{0.4cm}

  Afonso Sant'Anna

  \vspace{0.2cm}

  {\small Orientador: Yoshiharu Kohayakawa \quad
   Coorientador: Guilherme O.\ Mota}

  \vspace{0.2cm}

  {\small Instituto de Matemática e Estatística -- USP}
\end{center}

\vspace{0.5cm}

\noindent\textbf{Nota.}
Este documento é um rascunho de trabalho para o doutorado de Afonso Sant'Anna.
Ele reúne definições, teoremas e resultados dos artigos da literatura relevante,
organizados para facilitar a escrita dos próprios resultados do aluno.
Seções marcadas com \emph{(RA)} referem-se ao tema de Ramsey Assimétrico/Canônico;
seções marcadas com \emph{(AR)} referem-se ao tema Anti-Ramsey.

\tableofcontents

\bigskip

% ============================================================
\section{Preliminares e Notação}
% ============================================================

Denotamos por $\Gnp$ o grafo aleatório binomial com $n$ vértices em que cada aresta
está presente independentemente com probabilidade $p = p(n)$.
Dizemos que $\Gnp$ satisfaz uma propriedade \emph{assintoticamente quase certamente}
(\whp) se a probabilidade de satisfazê-la tende a $1$ quando $n \to \infty$.

\begin{definition}[Limiar e tipos de limiar]
Seja $\mathcal{P}$ uma propriedade monótona crescente de grafos.
Uma função $\hat{p} = \hat{p}(n)$ é um \textbf{limiar grosseiro} (\emph{coarse threshold})
para $\mathcal{P}$ em $\Gnp$ se
\[
  \lim_{n\to\infty} \Pr[\Gnp \in \mathcal{P}] =
  \begin{cases}
    0 & \text{se } p = o(\hat{p}),\\
    1 & \text{se } p = \omega(\hat{p}).
  \end{cases}
\]
Dizemos que $\hat{p}$ é um \textbf{limiar semi-agudo} (\emph{semi-sharp threshold})
se existem constantes $c, C > 0$ tais que
\[
  \lim_{n\to\infty} \Pr[\Gnp \in \mathcal{P}] =
  \begin{cases}
    0 & \text{se } p \le c\,\hat{p},\\
    1 & \text{se } p \ge C\,\hat{p}.
  \end{cases}
\]
O \textbf{$0$-enunciado} refere-se à afirmação para $p$ abaixo do limiar,
e o \textbf{$1$-enunciado} à afirmação para $p$ acima do limiar.
\end{definition}

\begin{definition}[Densidades de um grafo]
\label{def:densidades}
Dado um grafo $H$, definem-se:
\begin{itemize}
  \item A \textbf{densidade} de $H$:
    $d(H) := \dfrac{e(H)}{v(H)}$, onde $e(H) = |E(H)|$ e $v(H) = |V(H)|$.

  \item A \textbf{densidade máxima} de $H$:
    \[
      \maxdens(H) := \max\!\left\{\frac{e(J)}{v(J)} : J \subseteq H,\; v(J) \ge 1\right\}.
    \]
    Dizemos que $H$ é \textbf{balanceado} se $\maxdens(H) = d(H)$, e
    \textbf{estritamente balanceado} se $\maxdens(H) > d(J)$ para todo
    $J \subsetneq H$ com $v(J) \ge 1$.

  \item A \textbf{2-densidade} de $H$ (para $v(H) \ge 3$):
    \[
      d_2(H) := \frac{e(H)-1}{v(H)-2}.
    \]

  \item A \textbf{2-densidade máxima} de $H$:
    \[
      \mdens(H) := \max\!\left\{d_2(J) : J \subseteq H,\; v(J) \ge 3\right\}
        \cup \left\{\tfrac{1}{2}\right\}.
    \]
    (A fração $\frac{1}{2}$ corresponde ao caso $H \cong K_2$.)
    Dizemos que $H$ é \textbf{$2$-balanceado} se $\mdens(H) = d_2(H)$, e
    \textbf{estritamente $2$-balanceado} se $\mdens(H) > d_2(J)$ para todo
    $J \subsetneq H$.
\end{itemize}
\end{definition}

\begin{example}
  Para o clique $K_k$ com $k \ge 3$:
  $\mdens(K_k) = \dfrac{\binom{k}{2}-1}{k-2} = \dfrac{k+1}{2}$.
  Para o ciclo $C_\ell$ com $\ell \ge 3$:
  $\mdens(C_\ell) = \dfrac{\ell-1}{\ell-2}$.
  Em particular,
  $\mdens(K_5) = 3$, $\mdens(C_5) = \frac{4}{3}$, $\mdens(C_4) = \frac{3}{2}$.
\end{example}

% ============================================================
\section{Propriedade de Ramsey em Grafos Aleatórios}
% ============================================================

\begin{definition}[Propriedade $r$-Ramsey]
Para grafos $G$ e $H$ e um inteiro $r \ge 2$, escrevemos
$G \ramsey{r} H$ para denotar a propriedade de que,
em toda coloração das arestas de $G$ com $r$ cores,
existe pelo menos uma cópia monocromática de $H$ em $G$.
\end{definition}

O seguinte resultado é o marco fundamental da teoria:

\begin{theorem}[Rödl--Ruciński \cite{rodl1993lower, rodl1994random, rodl1995threshold}]
\label{thm:rodl-rucinski}
Seja $r \ge 2$ um inteiro e seja $H$ um grafo que não é uma floresta de estrelas
(e, no caso $r = 2$, também não é um caminho de comprimento~$3$).
Existem constantes positivas $c = c(H,r)$ e $C = C(H,r)$ tais que
\[
  \lim_{n\to\infty} \Pr\!\bigl[\Gnp \ramsey{r} H\bigr] =
  \begin{cases}
    0 & \text{se } p \le c\, n^{-1/\mdens(H)},\\
    1 & \text{se } p \ge C\, n^{-1/\mdens(H)}.
  \end{cases}
\]
Em particular, $\hat{p}(n) = n^{-1/\mdens(H)}$ é o limiar semi-agudo para
a propriedade $r$-Ramsey de $H$ em $\Gnp$.
\end{theorem}

Uma prova curta e elegante do Teorema~\ref{thm:rodl-rucinski} foi dada por
Nenadov e Steger \cite{nenadov2016short}, usando a técnica de
\emph{containers} de hipergrafos.

% ============================================================
\section{Propriedade Anti-Ramsey em Grafos Aleatórios} \label{sec:antiramsey}
% ============================================================

\begin{definition}[Propriedade anti-Ramsey]
\label{def:antiramsey}
Dados grafos $G$ e $H$, escrevemos $G \rb H$ para denotar a
\textbf{propriedade anti-Ramsey} de $G$ com respeito a $H$:
para toda coloração própria das arestas de $G$
(com número arbitrário de cores),
existe pelo menos uma cópia \emph{arco-íris} de $H$ em $G$,
isto é, uma cópia de $H$ em que cada aresta recebe uma cor distinta.
\end{definition}

\begin{remark}
A propriedade $G \rb H$ é monótona crescente em $G$.
Logo, para qualquer grafo fixo $H$, o grafo aleatório $\Gnp$ admite
um limiar $p_H^{\mathrm{rb}} = p_H^{\mathrm{rb}}(n)$ para essa propriedade
(cf.~Bollobás--Thomason).
\end{remark}

A questão central é determinar $p_H^{\mathrm{rb}}$ em termos de parâmetros
combinatórios de $H$.

\subsection{1-Enunciado: limite superior para o limiar}

O primeiro resultado geral foi obtido por Rödl e Tuza~\cite{rodl1992rainbow}
para ciclos, e depois estendido por Kohayakawa, Konstadinidis e Mota:

\begin{theorem}[Kohayakawa--Konstadinidis--Mota \cite{kohayakawa2014anti}]
\label{thm:1statement}
Seja $H$ um grafo qualquer. Existe uma constante $C > 0$ tal que,
se $p = p(n) \ge C\, n^{-1/\mdens(H)}$, então \whp\
$\Gnp \rb H$.
Em particular, $p_H^{\mathrm{rb}} \le n^{-1/\mdens(H)}$.
\end{theorem}

Em outras palavras, o limiar anti-Ramsey nunca está acima do limiar Ramsey clássico.

\subsection{0-Enunciado: limite inferior — o \emph{framework} de Nenadov et al.}

Para o limite inferior, Nenadov, Person, Škorić e Steger desenvolveram um
arcabouço geral que reduz o problema probabilístico a uma questão determinística:

\begin{theorem}[Nenadov--Person--Škorić--Steger \cite{nenadov2017algorithmic}]
\label{thm:framework}
Seja $H$ um grafo estritamente $2$-balanceado.
Se todo grafo $G$ tal que $G \rb H$ satisfaz $\maxdens(G) > \mdens(H)$,
então existe uma constante $c > 0$ tal que, se $p \le c\, n^{-1/\mdens(H)}$,
então \whp\ $\Gnp \not\rb H$.

Em outras palavras, verificar a condição determinística acima
é suficiente para obter o $0$-enunciado semi-agudo.
\end{theorem}

\begin{definition}[Condição determinística anti-Ramsey]
\label{def:cond-det}
Dado um grafo $H$, dizemos que $H$ \textbf{satisfaz a condição}~(D)
se todo grafo $G$ tal que $G \rb H$ satisfaz $\maxdens(G) > \mdens(H)$.
\end{definition}

\subsection{Resultados para classes específicas de grafos}

\subsubsection*{Ciclos}

\begin{theorem}[Barros--Cavalar--Mota--Parczyk \cite{AntiCycles}]
\label{thm:anticycles}
O limiar $p_{C_\ell}^{\mathrm{rb}}$ para a propriedade anti-Ramsey do ciclo $C_\ell$
satisfaz:
\begin{itemize}
  \item Se $\ell \ge 5$, então
    $p_{C_\ell}^{\mathrm{rb}} = n^{-1/\mdens(C_\ell)} = n^{-(\ell-2)/(\ell-1)}$
    (limiar semi-agudo).
  \item Para $\ell = 4$:
    $p_{C_4}^{\mathrm{rb}} = n^{-3/4} \ll n^{-1/\mdens(C_4)} = n^{-2/3}$.
  \item Para $\ell = 3$: a condição de ser grafo arco-íris é satisfeita
    trivialmente, logo
    $p_{C_3}^{\mathrm{rb}} = n^{-1}$ (limiar para a presença de um triângulo).
\end{itemize}
\end{theorem}

\subsubsection*{Grafos completos}

\begin{theorem}[Kohayakawa--Mota--Parczyk--Schnitzer \cite{AntiComplete}]
\label{thm:anticomplete}
Para o clique $K_k$, o limiar $p_{K_k}^{\mathrm{rb}}$ satisfaz:
\begin{itemize}
  \item Se $k \ge 5$, então
    $p_{K_k}^{\mathrm{rb}} = n^{-1/\mdens(K_k)} = n^{-2/(k+1)}$
    (limiar semi-agudo).
  \item Para $k = 4$:
    $p_{K_4}^{\mathrm{rb}} = n^{-7/15} \ll n^{-1/\mdens(K_4)} = n^{-2/5}$.
  \item Para $k = 3$: trivialmente $p_{K_3}^{\mathrm{rb}} = n^{-1}$.
\end{itemize}
\end{theorem}

\subsubsection*{Caso geral: 2-densidade alta}

\begin{theorem}[Kuperwasser \cite{kuperwasser2026anti}]
\label{thm:kuperwasser}
Seja $H$ um grafo com $\mdens(H) \ge 19$.
Existem constantes $c_0 < c_1$ tais que
\[
  \lim_{n\to\infty}\Pr\!\bigl[\Gnp \rb H\bigr] =
  \begin{cases}
    0 & \text{se } p \le c_0\, n^{-1/\mdens(H)},\\
    1 & \text{se } p \ge c_1\, n^{-1/\mdens(H)}.
  \end{cases}
\]
\end{theorem}

O Teorema~\ref{thm:kuperwasser} é provado verificando a condição determinística
do Teorema~\ref{thm:framework} através do seguinte resultado:

\begin{theorem}[Kuperwasser \cite{kuperwasser2026anti}]
\label{thm:kuperwasser-det}
Todo grafo $G$ com $\maxdens(G) \ge 18$ admite uma coloração própria das arestas
tal que todo subgrafo arco-íris tem $2$-densidade estritamente menor que $\maxdens(G)$.
Em particular, se $\mdens(H) > 18$, então $H$ satisfaz a condição~\emph{(D)}.
\end{theorem}

\subsection{A tabela de limiares conhecidos}

A tabela a seguir resume o estado atual dos limiares para a propriedade anti-Ramsey,
adaptada de \cite{behague2025thresholds}:

\begin{center}
\renewcommand{\arraystretch}{1.4}
\begin{tabular}{lll}
\hline
\textbf{Grafo $H$} & \textbf{Limiar $p_H^{\mathrm{rb}}$} & \textbf{Referência}\\
\hline
$K_\ell$, $\ell \ge 5$ & $n^{-2/(\ell+1)}$ (semi-agudo) & \cite{AntiComplete}\\
$K_4$                  & $n^{-7/15} \ll n^{-1/m_2(K_4)}$& \cite{AntiComplete}\\
$K_3$                  & $n^{-1}$                        & trivial\\
$C_\ell$, $\ell \ge 5$ & $n^{-(\ell-2)/(\ell-1)}$ (semi-agudo) & \cite{AntiCycles}\\
$C_4$                  & $n^{-3/4} \ll n^{-1/m_2(C_4)}$ & \cite{AntiCycles}\\
$C_3$                  & $n^{-1}$                        & trivial\\
$H$ com $m_2(H)\ge 19$ & $n^{-1/m_2(H)}$ (semi-agudo)   & \cite{kuperwasser2026anti}\\
\hline
\end{tabular}
\end{center}

\subsection{Questão em aberto: grafos de cintura grande}

Os resultados acima revelam um padrão: o limiar anti-Ramsey diverge de
$n^{-1/\mdens(H)}$ exatamente para grafos ``pequenos'' de cintura baixa
($K_3$, $K_4$, $C_4$). Essa observação motiva a seguinte questão central
deste projeto de doutorado:

\begin{question}
\label{q:cintura}
Seja $H$ um grafo com cintura (comprimento do menor ciclo) estritamente maior que~$4$.
É verdade que $p_H^{\mathrm{rb}} = n^{-1/\mdens(H)}$?
Equivalentemente, $H$ sempre satisfaz a condição~\emph{(D)}
(Definição~\ref{def:cond-det}) quando $\mathrm{girth}(H) > 4$?
\end{question}

% ============================================================
\section{Propriedade Ramsey Restrita (Constrained Ramsey)} \label{sec:cram}
% ============================================================

Este problema, estudado em \cite{behague2025thresholds}, generaliza tanto
a propriedade de Ramsey clássica quanto a anti-Ramsey.

\begin{definition}[Propriedade Ramsey restrita]
\label{def:cram}
Dados grafos $H_1$ e $H_2$, escrevemos $G \cram (H_1, H_2)$ para denotar
a \textbf{propriedade Ramsey restrita} (\emph{constrained Ramsey}):
para toda coloração própria das arestas de $G$,
existe uma cópia monocromática de $H_1$ \emph{ou} uma cópia arco-íris de $H_2$.
\end{definition}

\begin{remark}
Quando $H_2$ é uma floresta, qualquer coloração própria tem trivialmente uma
cópia arco-íris de $H_2$ (desde que o grafo seja grande o suficiente). Logo,
o caso interessante é quando $H_2$ não é uma floresta.

Quando $H_1 = $ estrela com $k$ arestas $K_{1,k}$ e $H_2 = H$, a propriedade
$G \cram (K_{1,k}, H)$ equivale a dizer que toda coloração \emph{localmente
$(k-1)$-limitada} de $G$ (em que cada cor aparece no máximo $k-1$ vezes por vértice)
contém uma cópia arco-íris de $H$. Assim,
a propriedade anti-Ramsey é o caso particular $k = 2$
(colorações próprias = colorações localmente $1$-limitadas).
\end{remark}

\subsection{Resultados principais de Behague et al.}

\begin{theorem}[Behague--Hancock--Hyde--Letzter--Morrison \cite{behague2025thresholds}]
\label{thm:behague-1-stmt}
Sejam $k \ge 2$ e $H$ um grafo qualquer.
Existe uma constante $C > 0$ tal que, se $p \ge C\, n^{-1/\mdens(H)}$, então
\whp\ $\Gnp \cram (K_{1,k}, H)$.
\end{theorem}

O Teorema~\ref{thm:behague-1-stmt} generaliza o Teorema~\ref{thm:1statement}.

\begin{definition}[Densidades relevantes para Ramsey restrito]
Para o problema restrito com $H_1 = K_{1,k}$, definem-se:
\begin{align*}
  m_a(H) &:= \inf\!\bigl\{\maxdens(G) : G \cram (K_{1,k}, H)\bigr\}\\
         &\phantom{:}= \text{densidade máxima do menor grafo} \ G\ \text{que força }\\
         &\phantom{= \text{densidade máxima do menor grafo} G \text{que força}} \text{Ramsey restrito para } (K_{1,k}, H).
\end{align*}
O limiar candidato natural é então $n^{-1/m_a(H)}$.
\end{definition}

\begin{theorem}[Behague et al.\ \cite{behague2025thresholds} -- $0$-enunciado geral]
\label{thm:behague-0-stmt}
Sejam $k \ge 3$ e $H$ um grafo que não é floresta.
\begin{enumerate}
  \item[(a)] Se $k \ge 4$, \emph{ou} $\mdens(H) \ne 2$, \emph{ou}
    $\mdens(H) = 2$ e $H$ contém um subgrafo estritamente $2$-balanceado
    $J \ne K_3$ com $\mdens(J) = 2$,
    então existem $c, C > 0$ tais que
    \[
      \lim_{n\to\infty} \Pr\!\bigl[\Gnp \cram (K_{1,k}, H)\bigr] =
      \begin{cases}
        0 & \text{se } p \le c\, n^{-1/\mdens(H)},\\
        1 & \text{se } p \ge C\, n^{-1/\mdens(H)}.
      \end{cases}
    \]
  \item[(b)] Se $k = 3$ e $\mdens(H) = 2$ e o único subgrafo estritamente
    $2$-balanceado de $H$ com $2$-densidade~$2$ é $K_3$ (e.g.\ $H = K_3$),
    então o limiar é mais fraco: coarse threshold em $n^{-1/2}$.
\end{enumerate}
\end{theorem}

% ============================================================
\section{Propriedade Ramsey Canônica Assimétrica} \label{sec:canonico}
% ============================================================

A teoria canônica de Ramsey estuda estruturas que emergem em colorações
\emph{arbitrárias} (sem restrição de número de cores).

\begin{definition}[Coloração canônica]
\label{def:canonica}
Seja $c : E(G) \to \mathbb{N}$ uma coloração arbitrária das arestas de $G$,
com uma ordenação total $v_1 < v_2 < \cdots < v_n$ dos vértices de $G$.
Um subgrafo $H \subseteq G$ possui uma \textbf{coloração canônica} se
satisfaz um dos seguintes padrões:
\begin{enumerate}
  \item \textbf{Monocromático:} todas as arestas de $H$ têm a mesma cor.
  \item \textbf{Arco-íris:} todas as arestas de $H$ têm cores distintas.
  \item \textbf{Lexicográfico (min-lex):} a cor de $v_iv_j$ (com $v_i < v_j$)
    depende unicamente de $v_i$.
  \item \textbf{Lexicográfico (max-lex):} a cor de $v_iv_j$ (com $v_i < v_j$)
    depende unicamente de $v_j$.
\end{enumerate}
\end{definition}
\noindent\hugo{definir a propriedade ramsey canônica}

\begin{theorem}[Erdős--Rado \cite{erdos1950combinatorial}]
Para qualquer inteiro $\ell \ge 1$, existe $N = N(\ell)$ tal que toda coloração
arbitrária das arestas de $K_N$ contém uma cópia canônica de $K_\ell$.
\end{theorem}

O resultado de Erdős--Rado foi trazido para o contexto aleatório por Kamčev e Schacht:

\begin{theorem}[Kamčev--Schacht \cite{kamcev2023canonical}]
\label{thm:kamcev}
Para todo $\ell \ge 4$, o limiar para a propriedade de que toda coloração
arbitrária de $\Gnp$ contenha uma cópia canônica de $K_\ell$ é
$\hat{p}(n) = n^{-2/(\ell+1)} = n^{-1/\mdens(K_\ell)}$.
\end{theorem}

\subsection{Cenário assimétrico}

\begin{definition}[Propriedade Ramsey canônica assimétrica]
\label{def:assimetrico}
Dados grafos $H_1$, $H_2$ e $H_3$, escrevemos
\[
  G \xrightarrow{\mathrm{can}} (H_1, H_2, H_3)
\]
para denotar a seguinte propriedade: para toda coloração arbitrária das
arestas de $G$, existe pelo menos
\begin{itemize}
  \item uma cópia \textbf{monocromática} de $H_1$, ou
  \item uma cópia \textbf{arco-íris} de $H_2$, ou
  \item uma cópia \textbf{lexicográfica} de $H_3$.
\end{itemize}
\end{definition}
\hugo{escolher qual definição de canônico usar. Todas as cópias lexicográficas de $H_3$, ou só uma cópia min-lex ou max-lex? Dizer que $G$ é ordenado.}

\begin{question}
\label{q:assimetrico}
Qual é o comportamento da função limiar $\hat{p}(n)$ para a propriedade
$\Gnp \xrightarrow{\mathrm{can}} (H_1, H_2, H_3)$
quando $H_1$, $H_2$ e $H_3$ não são todos iguais?
Em particular, como o limiar depende das densidades $\mdens(H_i)$?
\end{question}

% ============================================================
\section{O Caso Assimétrico \texorpdfstring{$(K_3, K_4, K_5)$}{(K3, K4, K5)}}
% ============================================================

Queremos determinar o limiar exato para a propriedade $\Gnp \xrightarrow{\mathrm{can}} (K_3,K_4,K_5)$.
Note que, para $p \ll n^{-1/\mdens(K_3)} = n^{-1/2}$, o grafo aleatório \whp\ não possui triângulos, tornando a propriedade trivialmente falsa (através de uma coloração arco-íris, por exemplo).
Por outro lado, para $p \gg n^{-1/\mdens(K_5)} = n^{-1/3}$, a existência de um $K_5$ canônico é garantida \whp, o que força a existência de uma das estruturas proibidas. \hugo{a.q.c. em vez de a.a.s.}

O resultado principal que buscamos formalizar é que o limiar coincide com o do clique mais denso.
Para a prova do $0$-enunciado, precisaremos das seguintes definições estruturais:

\begin{definition}[$H$-fechamento]
\label{def:H-fechado}
Dado um grafo $G$ e um grafo $H$, dizemos que:
\begin{itemize}
  \item Uma aresta $e \in E(G)$ é \textbf{$H$-fechada} se pertence a pelo menos duas cópias de~$H$ em~$G$.
  \item Uma cópia $H' \subseteq G$ de $H$ é \textbf{$H$-fechada} se pelo menos três de suas arestas são $H$-fechadas.
  \item O grafo $G$ é \textbf{$H$-fechado} se todo vértice e toda aresta de~$G$ pertencem a pelo menos uma cópia de~$H$, e toda cópia de~$H$ em~$G$ é $H$-fechada.
\end{itemize}
\end{definition}

\begin{definition}[$H$-Bloco]
\label{def:H-bloco}
Um grafo $G$ é um \textbf{$H$-bloco} se $G$ é $H$-fechado e, para todo subconjunto próprio não-vazio de arestas $E' \subsetneq E(G)$, existe uma cópia $H'$ de $H$ em $G$ tal que $E(H') \cap E' \ne \emptyset$ e $E(H') \setminus E' \ne \emptyset$.
\end{definition}

\begin{remark}
A chave do framework de Nenadov et al.\ é que, quando $p \ll n^{-1/\maxdens(H)}$\hugo{$m_2(H)$?}, os $H$-blocos presentes em $\Gnp$ têm tamanho $O(1)$ \whp.
\end{remark}

\hugo{enunciar o teorema entes das definições estruturais.}
\begin{theorem}
  Existe um limiar semi-agudo para a propriedade $\Gnp \xrightarrow{\mathrm{can}} (K_3,K_4,K_5)$ em $\hat{p}(n) = n^{-1/3}$. Ou seja, existem constantes $c, C > 0$ tais que:
  \[
  \lim_{n \to \infty} \mathbb{P}(\Gnp \xrightarrow{\mathrm{can}} (K_3,K_4,K_5)) = 
  \begin{cases}
    1 & \text{se } p \ge C n^{-1/3}, \\
    0 & \text{se } p \le c n^{-1/3}.  
  \end{cases}
  \]
\end{theorem}

\begin{proof}[Esboço da Prova]
  O $1$-enunciado segue do Teorema de Kam\v{c}ev--Schacht~\cite{kamcev2023canonical}.

  Para o $0$-enunciado, descrevemos o algoritmo \textsc{ColoraçãoCan} que, dado $G$ com $\maxdens(G) < 3$, constrói explicitamente uma coloração das arestas de $G$ sem $K_3$ monocromático, $K_4$ arco-íris ou $K_5$ lexicográfico.

  \textbf{Passo 1 --- Coloração Lexicográfica Inicial.}
  Fixamos uma ordenação total $\sigma: V(G) \to [n]$ e uma atribuição $\varphi: V(G) \to \mathbb{N}$.
  Definimos a coloração min-lex $\chi: E(G) \to \mathbb{N}$ por
  \[ \chi(uv) \;:=\; \varphi\!\bigl(\min_\sigma\{u,v\}\bigr), \]
  ou seja, a cor de cada aresta é determinada pelo vértice de menor índice segundo $\sigma$.
  Por construção, $\chi$ é uma coloração min-lex e satisfaz:
  \begin{itemize}
    \item não existe $K_3$ monocromático sob $\chi$;
    \item não existe $K_4$ arco-íris sob $\chi$.
  \end{itemize}
  O problema é que toda cópia de $K_5$ em $G$ é lexicográfica (min-lex) sob $\chi$ --- exatamente a estrutura proibida que precisamos eliminar.

  \textbf{Passo 2 --- Remoção de Arestas Livres.}
  Removemos de $G$ todas as arestas que não pertencem a nenhuma cópia de $K_5$.
  Seja $G' \subseteq G$ o grafo resultante. As arestas removidas mantêm a coloração $\chi$ do Passo~1 e não causam problemas.

  \textbf{Passo 3 --- Decomposição em $K_5$-Blocos.}
  Decompomos $G'$ em seus $K_5$-blocos (Definição~\ref{def:H-bloco}).
  Como $\maxdens(G) < \mdens(K_5) = 3$, cada $K_5$-bloco tem tamanho $O(1)$.

  \textbf{Passo 4 --- Recoloração dos Blocos.}
  Para cada $K_5$-bloco $B$, substituímos $\chi|_{E(B)}$ por uma nova coloração $\chi'$ (construída indutivamente) que evita $K_5$ lexicográfico, sem introduzir $K_3$ monocromático ou $K_4$ arco-íris.
  A existência de tal coloração em cada bloco é garantida pelo Lema~\ref{lem:0-statement-K5} abaixo.
\end{proof}

\begin{lemma}
  \label{lem:0-statement-K5}
  Seja $B$ um $K_5$-bloco com $\maxdens(B) < 3$. Então existe uma coloração $c: E(B) \to \mathbb{N}$ tal que:
  \begin{enumerate}
    \item[(i)] $c$ não possui $K_3$ monocromático;
    \item[(ii)] $c$ não possui $K_4$ arco-íris;
    \item[(iii)] $c$ não possui $K_5$ lexicográfico.
  \end{enumerate}
\end{lemma}

\begin{proof}

\noindent\textit{A coloração auxiliar de $K_5$.}

Seja $B$ um $K_5$-bloco com $\maxdens(B) < 3$.
Consideramos uma ordenação $\sigma$ dos vértices de $B$ e a coloração lexicográfica min-lex
$\chi(uv) := \varphi(\min_\sigma\{u,v\})$.

O ponto de partida é uma única cópia $H_0 \cong K_5$ com vértices ordenados $v_1 < v_2 < v_3 < v_4 < v_5$.
Sob a coloração min-lex, as cores das arestas são distribuídas como segue:

\medskip
\medskip
\begin{center}
\begin{tikzpicture}[scale=1.3,
  vx/.style={circle,draw,fill=white,minimum size=16pt,font=\small}]
  % K5 vertices as regular pentagon
  \node[vx] (1) at (90:1.5)  {$v_1$};
  \node[vx] (2) at (18:1.5)  {$v_2$};
  \node[vx] (3) at (-54:1.5) {$v_3$};
  \node[vx] (4) at (-126:1.5){$v_4$};
  \node[vx] (5) at (162:1.5) {$v_5$};
  % color a = phi(v1): edges 12,13,14,15
  \draw[line width=1.8pt, red]    (1)--(2);
  \draw[line width=1.8pt, red]    (1)--(3);
  \draw[line width=1.8pt, red]    (1)--(4);
  \draw[line width=1.8pt, red]    (1)--(5);
  % color b = phi(v2): edges 23,24,25
  \draw[line width=1.8pt, blue]   (2)--(3);
  \draw[line width=1.8pt, blue]   (2)--(4);
  \draw[line width=1.8pt, blue]   (2)--(5);
  % color c = phi(v3): edges 34,35
  \draw[line width=1.8pt, teal]   (3)--(4);
  \draw[line width=1.8pt, teal]   (3)--(5);
  % color d = phi(v4): edge 45
  \draw[line width=1.8pt, orange] (4)--(5);
\end{tikzpicture}
\end{center}

\noindent
É fácil verificar que toda cópia de $K_5$ sob $\chi$ é lexicográfica (min-lex).

Definimos agora uma \textbf{coloração auxiliar de duas cores} $\chi': E(B) \to \{\text{claro}, \text{escuro}\}$
(representadas como vermelho e azul nos diagramas) que servirá de base para a recoloração dos $K_5$-blocos.
Sobre o $K_5$ inicial $H_0$, $\chi'$ é definida identificando os $10$ arestas de $K_5$ com as $10$ arestas de
duas cópias de $C_5$: o \textbf{pentágono externo} (as arestas $v_1v_2, v_2v_3, v_3v_4, v_4v_5, v_5v_1$)
e o \textbf{pentágono interno} (as diagonais $v_1v_3, v_3v_5, v_5v_2, v_2v_4, v_4v_1$).

\begin{center}
\begin{tikzpicture}[scale=1.3,
  vx/.style={circle,draw,fill=white,minimum size=16pt,font=\small}]
  \node[vx] (1) at (90:1.5)  {$v_1$};
  \node[vx] (2) at (18:1.5)  {$v_2$};
  \node[vx] (3) at (-54:1.5) {$v_3$};
  \node[vx] (4) at (-126:1.5){$v_4$};
  \node[vx] (5) at (162:1.5) {$v_5$};
  % outer C5 (pentagon) -- claro = vermelho
  \draw[line width=2pt, red]  (1)--(2)--(3)--(4)--(5)--(1);
  % inner C5 (pentagram) -- escuro = azul
  \draw[line width=2pt, blue] (1)--(3)--(5)--(2)--(4)--(1);
\end{tikzpicture}
\end{center}

\noindent
Verifica-se diretamente que $\chi|_{H_0}$ é uma coloração própria de $H_0$ (todo vértice tem grau $2$ em cada $C_5$)\hugo{tá usando qual definição de coloração própria?},
não possui $K_3$ monocromático, não possui $K_4$ arco-íris, e não possui $K_5$ lexicográfico.


\medskip
\noindent\textit{Função de penalidade e tipos de expansão.}

Para controlar a estrutura dos $K_5$-blocos durante a construção indutiva,
introduzimos a seguinte medida adaptada de~\cite{AntiComplete}:

\begin{definition}[Função de penalidade]
Dado um grafo $G$, definimos
\[
  p(G) \;:=\; e(G) - 3\,v(G) + 5.
\]
\end{definition}

A função $p$ é normalizada de forma que o $K_5$ inicial tenha penalidade nula:
\[
  p(K_5) = 10 - 3 \cdot 5 + 5 = 0.
\]
Além disso, a condição de densidade $\maxdens(G) < 3$ garante que $e(G) < 3\,v(G)$, o que implica
$e(G) \le 3\,v(G) - 1$. Portanto, $p(G) \le 4$ para qualquer subgrafo de $B$. Essa limitação superior é crucial para controlar a 
complexidade estrutural do bloco.

A construção de um $K_5$-bloco $B$ pode ser vista como um processo indutivo. Começamos com uma 
cópia base $H_0 \cong K_5$ e, a cada passo, anexamos novos vértices e arestas para formar novas 
cópias de $K_5$. Seguindo as estratégias desenvolvidas em~\cite{nenadov2017algorithmic} e \cite{AntiComplete},
classificamos essas extensões em quatro tipos estruturais (ver analogia aos casos de expansão):

\begin{itemize}
  \item \textbf{$R_\ell$ (Expansão regular):} Adicionam-se $3$ novos vértices e $9$ novas arestas que, juntamente 
    com uma aresta já existente, formam um novo $K_5$. 
    Neste caso, a variação da penalidade é $\Delta p = 9 - 3(3) = 0$, ou seja, a penalidade do 
    grafo se mantém inalterada.

  \item \textbf{$T_\ell$ (Expansão sobre triângulo):} Adicionam-se $2$ novos vértices e $7$ novas arestas que, 
    apoiadas em um triângulo preexistente, completam um novo $K_5$. 
    Aqui, o incremento na penalidade é $\Delta p = 7 - 3(2) = 1$.

  \item \textbf{$Q_\ell$ (Expansão sobre $K_4$):} Adiciona-se apenas $1$ novo vértice ligado a um $K_4$ 
    existente através de $4$ novas arestas, completando um $K_5$. 
    O incremento é $\Delta p = 4 - 3(1) = 1$.

  \item \textbf{$U_\ell$ (União sobre $K_5^-$):} Não se adicionam vértices; apenas $1$ nova aresta é 
    inserida conectando dois vértices de um $K_5$ quase completo ($K_5 \setminus \{e\}$) já 
    presente na estrutura. 
    A variação é $\Delta p = 1 - 3(0) = 1$.
\end{itemize}

Uma observação fundamental é que a expansão \textbf{$R_\ell$} (a única que é considerada "regular") 
preserva o valor de $p(G)$. Por outro lado, as expansões degeneradas (\textbf{$T_\ell$}, \textbf{$Q_\ell$} e \textbf{$U_\ell$}) 
aumentam a penalidade estritamente em $1$. Como iniciamos com $p(H_0) = 0$ e sabemos que $p(B) \le 4$, concluímos 
que o número total de ocorrências dos casos $T_\ell$, $Q_\ell$ e $U_\ell$ na construção do bloco é uniformemente 
limitado (ocorrem no máximo $4$ vezes em todo o processo).

\medskip
\noindent\textit{Extensão da coloração auxiliar para cada tipo de expansão.}
\medskip

\noindent\textbf{Caso $R_\ell$ (expansão regular).}
Em $R_\ell$, uma nova cópia de $K_5$ é colada à estrutura existente por uma aresta $e$ já colorida.
Identificamos os dois $C_5$ na nova cópia de modo que $e$ pertença ao $C_5$ da sua cor.
A coloracão é compatível com o bloco anterior.
No diagrama abaixo, dois $K_5$ são colados pela aresta $1$--$2$ (vermelha):

\medskip
\begin{center}
\begin{tikzpicture}[scale=1.0,
  vx/.style={circle,draw,fill=white,minimum size=13pt,font=\tiny\strut}]
  \node[vx] (v1) at (0, 0.9)   {$1$};
  \node[vx] (v2) at (0,-0.9)   {$2$};
  \node[vx] (v5) at (-1.4, 1.6){$5$};
  \node[vx] (v4) at (-2.0, 0)  {$4$};
  \node[vx] (v3) at (-1.4,-1.6){$3$};
  \node[vx] (v8) at ( 1.4, 1.6){$8$};
  \node[vx] (v7) at ( 2.0, 0)  {$7$};
  \node[vx] (v6) at ( 1.4,-1.6){$6$};
  % Left K5: Red outer C5: 1-5-4-3-2-1
  \draw[red, line width=1.8pt]  (v1)--(v5)--(v4)--(v3)--(v2);
  % Left K5: Blue inner C5: 1-4-2-5-3-1
  \draw[blue, line width=1.8pt] (v1)--(v4) (v4)--(v2) (v2)--(v5) (v5)--(v3) (v3)--(v1);
  % Right K5: Red outer C5: 1-8-7-6-2-1
  \draw[red, line width=1.8pt]  (v1)--(v8)--(v7)--(v6)--(v2);
  % Right K5: Blue inner C5: 1-7-2-8-6-1
  \draw[blue, line width=1.8pt] (v1)--(v7) (v7)--(v2) (v2)--(v8) (v8)--(v6) (v6)--(v1);
  % Shared edge highlighted
  \draw[red, line width=3pt] (v1)--(v2);
  \node[font=\tiny] at (0.25,0) {$e$};
\end{tikzpicture}
\end{center}

\medskip
\noindent\textbf{Caso $T_\ell$ (expansão sobre triângulo).}
Em $T_\ell$, dois novos vértices $u, v$ com a aresta $uv$ são ligados a um triângulo $\{a,b,c\}$ existente.
Identificamos um $P_2$ monocromático no triângulo (sempre existe por paridade) e extendemos
ao $C_5$ da sua cor; as cinco arestas restantes formam o $C_5$ da outra cor.
Se o $P_2$ vermelho é $a$--$b$--$c$, então: $C_5$ vermelho $= a\text{-}b\text{-}c\text{-}y\text{-}x\text{-}a$;
$C_5$ azul $= a\text{-}c\text{-}x\text{-}b\text{-}y\text{-}a$.

\medskip
\begin{center}
\begin{tikzpicture}[scale=1,
  vx/.style={circle,draw,fill=white,minimum size=13pt,font=\tiny\strut},
  ov/.style={circle,draw,fill=gray!25,minimum size=13pt,font=\tiny\strut}]

%================================================
% K5 com coordenadas explícitas
%================================================

\node[ov] (A) at (0,1.8)    {$a$};
\node[ov] (B) at (1.7,0.6)  {$b$};
\node[ov] (C) at (1.0,-1.4) {$c$};
\node[vx] (U) at (-1.0,-1.4) {$u$};
\node[vx] (V) at (-1.7,0.6) {$v$};

% vermelho: a-b-c-u-v-a
\draw[red, line width=1.8pt]
(A)--(B)--(C)--(U)--(V)--(A);

% azul: a-c-v-b-u-a
\draw[blue, line width=1.8pt]
(A)--(C)
(C)--(V)
(V)--(B)
(B)--(U)
(U)--(A);

%================================================
% vértices extras x,y à direita
%================================================

\node[vx] (X) at (4.0,0.8)  {$x$};
\node[vx] (Y) at (4.0,-0.8) {$y$};

% aresta xy
\draw[black, line width=1.5pt] (X)--(Y);

% ligações ao triângulo abc
\draw[dashed,gray]
(X) to[bend right=20](A)
(X)--(B)
(X)--(C);

\draw[dashed,gray]
(Y) to[bend right=20](A)
(Y) to[bend left=20] (B)
(Y)--(C);

\end{tikzpicture}
\end{center}
\medskip
\noindent\textbf{Caso $Q_\ell$ (expansão sobre $K_4$).}
Em $Q_\ell$, um novo vértice $v$ é ligado a todos os vértices de um $K_4$ existente $\{a,b,c,d\}$.
O $K_4$ possui dois $P_3$ monocromáticos; tomamos o $P_3$ vermelho $a$--$b$--$c$--$d$ e extendemos:
$C_5$ vermelho $= a\text{-}b\text{-}c\text{-}d\text{-}v\text{-}a$;
$C_5$ azul $= a\text{-}c\text{-}v\text{-}b\text{-}d\text{-}a$.

\medskip
\begin{center}
\begin{tikzpicture}[scale=1,
  vx/.style={circle,draw,fill=white,minimum size=13pt,font=\tiny\strut},
  ov/.style={circle,draw,fill=gray!25,minimum size=13pt,font=\tiny\strut}]

%================================================
% K5
%================================================

\node[ov] (A) at (0,1.8)     {$a$};
\node[ov] (B) at (1.7,0.6)   {$b$};
\node[ov] (C) at (1.0,-1.4)  {$c$};
\node[vx] (U) at (-1.0,-1.4) {$u$};
\node[ov] (V) at (-1.7,0.6)  {$v$};

% vermelho: a-b-c-u-v-a
\draw[red, line width=1.8pt]
(A)--(B)--(C)--(U)--(V)--(A);

% azul: a-c-v-b-u-a
\draw[blue, line width=1.8pt]
(A)--(C)
(C)--(V)
(V)--(B)
(B)--(U)
(U)--(A);

%================================================
% vértice extra à direita
%================================================

\node[vx] (X) at (4.2,0) {$x$};

% x ligado ao K4 abcv
\draw[dashed,gray]
(X) to[bend right=20] (A)
(X)--(B)
(X)--(C)
(X) to[bend left=20] (V);

\end{tikzpicture}
\end{center}
\medskip

% ------------------------------------------------------------------
% ESTRUTURA DA PROVA: enunciados e deduções
% ------------------------------------------------------------------

\begin{proposition}[Equilíbrio de cores]
\label{prop:equilibrio}
Seja $B$ um $K_5$-bloco com $\maxdens(B) < 3$ e $\chi'$ a coloração auxiliar construída
indutivamente sobre as expansões de $B$. Então é possível definir $\chi'$ de modo que,
para toda cópia $H \cong K_5$ em $B$ e todo vértice $w \in V(H)$, $w$ possui ao menos
um vizinho com $\chi'(e) = \text{claro}$ e ao menos um vizinho com $\chi'(e) = \text{escuro}$
dentro de $H$.
\end{proposition}

\noindent
Para provar a Proposição~\ref{prop:equilibrio}, estabelecemos duas afirmações auxiliares.

\begin{afirmacao}[Expansões de vértice preservam o equilíbrio]
\label{afirm:vertex-exp}
A propriedade da Proposição~\ref{prop:equilibrio} é preservada em cada passo dos tipos
$R_\ell$, $T_\ell$ e $Q_\ell$.
\end{afirmacao}

\begin{afirmacao}[Limitação de $U_\ell$]
\label{afirm:ulimit}
Na construção de qualquer $K_5$-bloco $B$ com $\maxdens(B) < 3$,
o número de expansões do tipo $U_\ell$ é no máximo $2$.
\end{afirmacao}

\begin{proof}[Prova da Proposição~\ref{prop:equilibrio}]
Pela Afirmação~\ref{afirm:vertex-exp}, a propriedade de equilíbrio é mantida a cada
expansão dos tipos $R_\ell$, $T_\ell$ e $Q_\ell$.
Pela Afirmação~\ref{afirm:ulimit}, o número de passos $U_\ell$ é limitado;
e cada passo $U_\ell$ apenas adiciona uma aresta, sem alterar a propriedade já garantida para
os três vizinhos existentes de cada extremidade (como detalhado na prova da Afirmação~\ref{afirm:ulimit}).
Portanto, $\chi'$ pode ser construída de modo que a propriedade vale após todos os passos.
\end{proof}

\begin{afirmacao}[$\chi''$ é uma coloração válida]
\label{afirm:chipp}
Dadas $\chi$ e $\chi'$ como acima, a coloração
\[
  \chi'': E(B) \;\longrightarrow\; \mathbb{N} \times \{\text{claro}, \text{escuro}\},
  \qquad \chi''(e) := \bigl(\chi(e),\, \chi'(e)\bigr),
\]
evita $K_3$ monocromático, $K_4$ arco-íris e $K_5$ lexicográfico.
Em particular, $\chi''$ é uma $(K_3, K_4, K_5)$-coloração válida de $B$,
o que conclui a prova do Lema~\ref{lem:0-statement-K5}.
\end{afirmacao}

\begin{proof}
Assumimos que $\chi'$ satisfaz a Proposição~\ref{prop:equilibrio} e verificamos cada condição
separadamente.

\begin{itemize}
  \item Sob $\chi$, toda cópia de $K_5$ é lexicográfica: as quatro arestas $v_1v_j$
  incidentes ao vértice mínimo $v_1$ recebem todas a mesma cor $\chi(v_1v_j) = \varphi(v_1)$.
  Para que o $K_5$ seja também min-lex sob $\chi''$, seria necessário que todas essas
  arestas tivessem o mesmo valor de $\chi'$. Isso é impossível pela
  Proposição~\ref{prop:equilibrio}, que garante que $v_1$ possui ao menos uma aresta
  clara e uma escura dentro de qualquer cópia de $K_5$. Portanto, não existe
  $K_5$ lexicográfico sob $\chi''$.

  \item Considere um $K_4$ com vértice mínimo $w$. As três arestas $wu_i$ incidentes
  a $w$ recebem todas a mesma primeira componente $\chi(wu_i) = \varphi(w)$; para que
  o $K_4$ fosse arco-íris sob $\chi''$, essas três arestas precisariam ser duas a duas
  distintas na segunda componente $\chi'$. Como só há duas cores em $\chi'$, ao menos
  duas delas coincidem em $\chi'$ e, portanto, em $\chi''$. Logo, não existe
  $K_4$ arco-íris sob $\chi''$.

  \item Em qualquer triângulo $x < y < z$, as arestas $xy$ e $xz$ recebem $\chi(xy) =
  \varphi(x)$, enquanto $yz$ recebe $\chi(yz) = \varphi(y) \ne \varphi(x)$ (pois $\varphi$
  é injetiva). Assim, $\chi''(xy) \ne \chi''(yz)$, e o triângulo não pode ser
  monocromático sob $\chi''$.
\end{itemize}
\end{proof}

\noindent
A Afirmação~\ref{afirm:chipp} conclui a prova do Lema~\ref{lem:0-statement-K5}.

\end{proof}

% ------------------------------------------------------------------
% PROVAS DAS AFIRMAÇÕES AUXILIARES
% ------------------------------------------------------------------

\begin{proof}[Prova da Afirmação~\ref{afirm:vertex-exp} (Expansões de vértice)]
A prova é por indução no número de passos de expansão na construção do bloco $B$.
Mantemos o seguinte invariante: após cada passo, toda cópia de $K_5$ presente no grafo
atual admite uma decomposição $E(K_5) = E(C_5^{\text{vm}}) \cup E(C_5^{\text{az}})$
(com todo vértice de grau $2$ em cada $C_5$), e todo triângulo presente tem exatamente
$2$ arestas de uma cor e $1$ da outra (distribuição $2$-$1$).

\medskip

O caso base $H_0 \cong K_5$ satisfaz o invariante pela construção explícita da coloração
auxiliar $\chi'$ descrita anteriormente.

\medskip\noindent\textbf{Passo indutivo --- Caso $R_\ell$.}
Uma nova cópia de $K_5$ é colada à estrutura existente por uma única aresta $e$ já colorida
(digamos, $e$ é vermelha). As $9$ novas arestas ainda não possuem cor.
Como $K_5$ se decompõe em dois $C_5$, e qualquer aresta de $K_5$ pode pertencer ao $C_5$
de qualquer uma das duas cores, basta escolher a decomposição do novo $K_5$ de modo que
$e$ pertença ao $C_5$ vermelho. Com isso, cada vértice do novo $K_5$ fica com grau $2$
em cada $C_5$, e o invariante é preservado. Além disso, como a aresta compartilhada $e$
está corretamente inserida no $C_5$ da sua cor, nenhum triângulo preexistente é afetado.
Note que esta extensão é sempre possível, independentemente da cor de $e$, e portanto
uma sequência arbitrariamente longa de passos $R_\ell$ preserva o invariante.

\medskip\noindent\textbf{Passo indutivo --- Caso $T_\ell$.}
Dois novos vértices $u, v$ (com a aresta $uv$) são ligados a um triângulo $\{a,b,c\}$
preexistente, formando um novo $K_5$ sobre $\{a,b,c,u,v\}$.
Pelo invariante indutivo, o triângulo $\{a,b,c\}$ possui distribuição $2$-$1$;
logo existe um $P_2$ monocromático, digamos o caminho vermelho $a$--$b$--$c$.
Definimos a coloração do novo $K_5$ estendendo esse $P_2$ ao $C_5$ vermelho
$a\text{-}b\text{-}c\text{-}u\text{-}v\text{-}a$; as cinco arestas restantes
$\{au, bu, bv, cv, uv\} \setminus \{av\}$
formam o $C_5$ azul $a\text{-}c\text{-}v\text{-}b\text{-}u\text{-}a$.
Verifica-se diretamente que todo vértice tem grau $2$ em cada $C_5$, preservando o invariante.

\medskip\noindent\textbf{Passo indutivo --- Caso $Q_\ell$.}
Um novo vértice $v$ é ligado a todos os vértices de um $K_4$ existente $\{a,b,c,d\}$,
formando um novo $K_5$ sobre $\{a,b,c,d,v\}$.
Pelo invariante, toda cópia de $K_5$ que contém o $K_4$ tem distribuição $2$-$1$ em seus
triângulos; em particular, o $K_4$ possui um $P_3$ monocromático, digamos o caminho
vermelho $a$--$b$--$c$--$d$.
Definimos a coloração do novo $K_5$ estendendo esse $P_3$ ao $C_5$ vermelho
$a\text{-}b\text{-}c\text{-}d\text{-}v\text{-}a$; as cinco arestas restantes formam
o $C_5$ azul $a\text{-}c\text{-}v\text{-}b\text{-}d\text{-}a$.
Todo vértice fica com grau $2$ em cada $C_5$, e o invariante é preservado.
\end{proof}

\begin{proof}[Prova da Afirmação~\ref{afirm:ulimit} (Limitação de $U_\ell$)]
Dividimos a prova em dois casos, conforme o bloco $B$ contenha ou não uma cópia de $K_6$.

\medskip
\noindent\textit{Caso 2: O bloco $B$ contém uma cópia de $K_6$.}

\medskip

Como a $2$-densidade máxima do bloco satisfaz $\maxdens(B) < 3$, o bloco $B$ não pode conter nenhuma cópia de $K_7$ como subgrafo (uma vez que $m(K_7) = 21/7 = 3$). Assim, a maior cópia de clique completo que $B$ pode conter é $K_6$.

Fixamos $H_0 \cong K_6$ como grafo inicial da construção indutiva de $B$. Como $B$ é um $K_5$-bloco e contém $K_6$, é possível iniciar a construção indutiva a partir deste subgrafo $K_6$ maximal e anexar os demais vértices por expansões sucessivas, seguindo a estrutura padrão descrita em~\cite{nenadov2017algorithmic} e \cite{AntiComplete}.

\medskip
\noindent\textit{Penalidade inicial.}
Calculamos diretamente a penalidade de $H_0 \cong K_6$:
\[
  p(K_6) \;=\; e(K_6) - 3\,v(K_6) + 5 \;=\; 15 - 3 \cdot 6 + 5 \;=\; 2.
\]
Como a densidade máxima do bloco satisfaz $\maxdens(B) < 3$, temos $e(B) \le 3\,v(B) - 1$, o que garante que $p(B) \le 4$. Logo, o orçamento total para passos degenerados a partir de $H_0$ é de no máximo $p(B) - p(K_6) \le 4 - 2 = 2$.

\medskip
\noindent\textit{Coloração auxiliar de $K_6$.}
Rotulamos os vértices de $K_6$ como $1, 2, 3, 4, 5, 6$ e definimos $\chi'$ por:
\begin{itemize}
  \item $C_6$ vermelho: $1\text{-}2\text{-}3\text{-}4\text{-}5\text{-}6\text{-}1$
        (as $6$ arestas do ciclo externo);
  \item $C_6$ azul: $1\text{-}3\text{-}5\text{-}1$ e $2\text{-}4\text{-}6\text{-}2$
        reunidos, ou seja, as $6$ arestas $13, 35, 15, 24, 46, 26$
        (o ciclo ``saltando de dois em dois'');
  \item $M_3$ azul: o emparelhamento perfeito $\{14, 25, 36\}$.
\end{itemize}
Tem-se $6 + 6 + 3 = 15 = \binom{6}{2}$ arestas, cobrindo exatamente $E(K_6)$.

\medskip

\begin{center}
\begin{tikzpicture}[scale=1.2,
  vm/.style={draw=red!70!black, very thick},
  az/.style={draw=blue!70!black, very thick, dashed},
  vert/.style={circle, fill=black, inner sep=1.8pt}]

  \foreach \i/\ang in {1/90, 2/30, 3/330, 4/270, 5/210, 6/150}
    \node[vert] (v\i) at (\ang:1.2cm) {};
  \foreach \i/\ang in {1/90, 2/30, 3/330, 4/270, 5/210, 6/150}
    \node at (\ang:1.55cm) {$\i$};

  % C6 vermelho (ciclo externo)
  \foreach \a/\b in {1/2, 2/3, 3/4, 4/5, 5/6, 6/1}
    \draw[vm] (v\a) -- (v\b);
  % M3 vermelho (diâmetros)
  \foreach \a/\b in {1/4, 2/5, 3/6}
    \draw[az] (v\a) -- (v\b);
  % C6 azul (saltos de 2)
  \foreach \a/\b in {1/3, 3/5, 5/1, 2/4, 4/6, 6/2}
    \draw[az] (v\a) -- (v\b);
\end{tikzpicture}
\end{center}

\medskip
\noindent\textit{Verificação das propriedades invariantes.}

A coloração acima satisfaz o invariante estabelecido na prova da
Afirmação~\ref{afirm:vertex-exp}:

\begin{itemize}
  \item Todo vértice de $K_6$ tem grau $2$ no $C_6$ vermelho, grau $2$ no $C_6$ azul
        e grau $1$ no $M_3$ vermelho, portanto grau total $3$ em arestas vermelhas e
        grau $2$ em arestas azuis.

  \item Qualquer cópia de $K_5$ em $K_6$ é obtida removendo um vértice $w$.
        Dentro desse $K_5$, cada vértice perde exatamente $1$ vizinho (o próprio $w$),
        mas mantém ao menos $1$ vizinho vermelho e $1$ vizinho azul, preservando o
        equilíbrio da Proposição~\ref{prop:equilibrio}.

  \item Todo triângulo de $K_6$ possui exatamente $2$ arestas vermelhas e $1$ azul,
        ou $1$ vermelha e $2$ azuis (pode-se verificar por inspeção direta nos
        $\binom{6}{3} = 20$ triângulos). Em particular, a distribuição $2$-$1$ vale,
        garantindo a aplicabilidade da Afirmação~\ref{afirm:vertex-exp}.
\end{itemize}

\medskip
\noindent\textit{Como podem surgir passos $U_\ell$ a partir de $K_6$.}

Um passo $U_\ell$ requer a presença de um $K_5^-$ no grafo atual, isto é, cinco
vértices faltando exatamente uma aresta para formar $K_5$.
No entanto, dentro de $K_6$ puro, qualquer conjunto de $5$ vértices induz um $K_5$
completo; logo \emph{não existe} $K_5^-$ dentro de $H_0 = K_6$.

Portanto, um passo $U_\ell$ nunca pode ser o primeiro passo após $H_0$: é
necessário realizar ao menos um passo de expansão degenerado do tipo $T_\ell$ ou
$Q_\ell$ antes, de modo a introduzir novos vértices e criar estruturas contendo um $K_5^-$.

Cada um desses passos degenerados ($T_\ell$ ou $Q_\ell$) consome exatamente $1$
unidade do orçamento de penalidade. Como o orçamento total disponível a partir de
$K_6$ é de no máximo $2$ (devido à densidade estrita $\maxdens(B) < 3$), esse primeiro
passo reduz o orçamento restante para no máximo $2 - 1 = 1$.

Visto que cada passo $U_\ell$ consome exatamente $1$ unidade de penalidade ($\Delta p = 1$),
o orçamento restante de no máximo $1$ permite realizar \emph{no máximo um} passo
$U_\ell$ em toda a construção. Portanto, o número total de passos $U_\ell$ quando
$H_0 \cong K_6$ é no máximo $1$, o que satisfaz o limite de no máximo $2$ exigido
pela Afirmação~\ref{afirm:ulimit}, concluindo o Caso~2.

\medskip
\noindent\textit{Caso 1: O bloco $B$ não contém nenhuma cópia de $K_6$.}

\medskip

Iniciamos a construção a partir de uma única cópia base $H_0 \cong K_5$, com
$p(H_0) = 0$. O orçamento total para passos degenerados é $p(B) \le 4$.

\medskip
\noindent\textit{Análise dos vértices acoplados via $Q_\ell$ a $H_0$.}

Na construção do bloco, a partir de $H_0$, podemos ter vértices adicionados diretamente sobre $H_0$ através de passos $Q_\ell$. Cada passo $Q_\ell$ adiciona um novo vértice conectado a um $K_4 \subset H_0$, consumindo $1$ unidade do orçamento de penalidade. Como o orçamento total é $4$, podemos ter até $4$ vértices acoplados dessa forma a $H_0$.

Seja $k$ o número de vértices acoplados a $H_0$ via passos $Q_\ell$. O orçamento consumido por esses passos iniciais é $k$, restando $4 - k$ de orçamento disponível para o resto da construção (incluindo passos $U_\ell$). Analisamos os casos em que $k \in \{2, 3, 4\}$:

\begin{itemize}
    \item \textbf{$k = 4$ vértices acoplados:} O orçamento consumido é $4$. Resta $0$ de orçamento. Portanto, o número de passos $U_\ell$ possíveis na construção é \textbf{$0$}.
    \item \textbf{$k = 3$ vértices acoplados:} O orçamento consumido é $3$. Resta apenas $1$ unidade de orçamento. Portanto, podemos ter no máximo \textbf{$1$} passo $U_\ell$ na construção.
    \item \textbf{$k = 2$ vértices acoplados:} O orçamento consumido é $2$. Restam $2$ unidades de orçamento. Como cada passo $U_\ell$ consome $1$, podemos ter no máximo \textbf{$2$} passos $U_\ell$.
\end{itemize}

Em todos esses casos ($k \in \{2, 3, 4\}$), é trivialmente óbvio que o número de passos $U_\ell$ é no máximo $2$, satisfazendo a Afirmação~\ref{afirm:ulimit}. O único cenário em que o orçamento restante permitiria, teoricamente, $3$ passos $U_\ell$ ocorre quando temos apenas \textbf{$1$ vértice acoplado} a $H_0$ via $Q_\ell$ (isto é, $k = 1$). Faremos uma análise detalhada deste caso.

\medskip
\noindent\textit{Passo $U_\ell$ imediatamente após um único $Q_\ell$ é impossível sem $K_6$.}

Suponha que, a partir de $H_0$, realizamos um único passo $Q_\ell$: adicionamos um vértice $u$
ligado a todos os vértices de um $K_4 = \{a,b,c,d\} \subset K_5$, formando um segundo
$K_5 = \{u,a,b,c,d\}$. Neste ponto $N(v) \cap N(u) = \{a,b,c,d\}$, onde $v$ é o vértice
original do primeiro $K_5$ fora de $K_4$. A vizinhança comum $\{a,b,c,d\}$ contém
um $K_4$, de modo que a aresta $vu$ completaria um $K_6$ sobre $\{v,u,a,b,c,d\}$. Como
$B$ não contém $K_6$, nenhum passo $U_\ell$ é possível imediatamente após $Q_\ell$.
É portanto necessário intercalar ao menos um passo $T_\ell$ antes do primeiro $U_\ell$.

\medskip
\noindent\textit{Configuração $H_1$ após $T_\ell$.}

Seja $H_0$ o grafo com dois $K_5$ compartilhando $K(v) = \{a,b,c,d\}$, obtido por um
passo $Q_\ell$ a partir do $K_5$ inicial (com $p(H_0) = 1$). Como $B$ não contém $K_6$,
realizamos um passo $T_\ell$: adicionamos dois vértices $x, y$ (com a aresta $xy$) ligados
a um triângulo $\Delta \subset B$. Consideramos dois subcasos conforme a posição de
$\Delta$ em relação a $K(v)$.

\medskip
\noindent\textbf{Subcaso 1a: $\Delta \subset K(v)$.}
Tomamos $\Delta = \{a,b,c\} \subset K(v)$. Após $T_\ell$, o grafo $H_1$ tem:
\[
    v(H_1) = 8,\quad e(H_1) = 21,\quad p(H_1) = 2.
\]
As únicas não-arestas de $H_1$ relevantes para $U_\ell$ são da forma $\{w, z\}$
com $w \in \{v, u\}$ e $z \in \{x, y\}$, ou $\{x,d\}$ e $\{y,d\}$: em todos os
casos, $N(w) \cap N(z) = \{a,b,c\}$, que contém um triângulo mas não um $K_4$
(pois $d$ não está ligado a $x$ nem a $y$), garantindo a ausência de $K_6$ ao
adicionar a aresta.

\medskip
\noindent\textbf{Subcaso 1b: $\Delta \not\subset K(v)$.}
O triângulo $\Delta$ contém pelo menos um vértice de $\{v,u\}$; tomamos sem perda
de generalidade $\Delta = \{v,a,b\}$. Após $T_\ell$:
\[
    v(H_1) = 8,\quad e(H_1) = 21,\quad p(H_1) = 2.
\]
As não-arestas relevantes para $U_\ell$ são $\{x,c\}$, $\{x,d\}$, $\{y,c\}$, $\{y,d\}$,
com $N(x) \cap N(z) = \{v,a,b\}$ para $z \in \{c,d\}$; as não-arestas $\{u,x\}$ e $\{u,y\}$
satisfazem $N(u) \cap N(x) = \{a,b\}$, que não contém triângulo, e portanto não dão
origem a $K_5^-$ nem a passos $U_\ell$.

\medskip
Em ambos os subcasos, a estrutura das não-arestas de $H_1$ tem a seguinte forma:
existe um triângulo base $\Delta = \{a,b,c\}$ e dois conjuntos de vértices
$A = \{v, u\}$ e $Z = \{x, y\}$ (no Subcaso~1a) ou $A = \{c,d\}$ e $Z = \{x,y\}$
(no Subcaso~1b), tais que:
\begin{itemize}
  \item todo $K_5^-$ disponível em $H_1$ tem a forma $\{w,z,a,b,c\}$ com
        $w \in A$, $z \in Z$ e aresta faltante $\{w,z\}$;
  \item $N(w) \cap N(z) = \{a,b,c\}$ para todo par $(w,z) \in A \times Z$.
\end{itemize}

\medskip
\noindent\textit{Ao mais dois passos $U_\ell$ sem criar $K_6$.}

Após $H_1$, os possíveis passos $U_\ell$ correspondem à adição de uma aresta
$\{w,z\} \in A \times Z = \{w_1,w_2\} \times \{z_1,z_2\}$. Analisamos quantos
desses passos são realizáveis sem criar $K_6$.

\medskip
Após o \textbf{primeiro passo} $U_\ell$, digamos a adição de $\{w_1, z_1\}$:
\begin{itemize}
  \item $N(w_1) \cap N(z_2) = \{a,b,c\}$: sem $K_4$, logo o segundo passo
        $\{w_1,z_2\}$ criaria $K_5$ sem $K_6$. \checkmark
  \item Analogamente para $\{w_2,z_1\}$.
  \item Para $\{w_2,z_2\}$: $N(w_2) \cap N(z_2)$ ainda não contém $w_1$ nem $z_1$
        (pois a aresta $w_2 z_2$ não impõe ligação com $w_1$ ou $z_1$),
        logo esse passo também não cria $K_6$. \checkmark
\end{itemize}
Portanto o \textbf{segundo passo} $U_\ell$ é sempre possível sem $K_6$.

\medskip
Após \textbf{dois passos} $U_\ell$, suponha que foram adicionadas as arestas
$\{w_i, z_j\}$ e $\{w_k, z_\ell\}$. Consideramos os dois casos:

\begin{enumerate}
  \item \emph{Par adjacente} (compartilhando um vértice de $A$ ou $Z$):
    os seis vértices $\{a,b,c,w_i,z_j,z_\ell\}$ (ou $\{a,b,c,w_i,w_k,z_j\}$)
    têm agora todas as arestas presentes, formando $K_6$. Logo este caso já violou
    a condição de ausência de $K_6$ ao ser executado. Assim, somente pares
    \emph{não-adjacentes} (disjuntos em $A$ e $Z$) são realizáveis: $\{w_1,z_1\}$
    e $\{w_2,z_2\}$, ou $\{w_1,z_2\}$ e $\{w_2,z_1\}$.

  \item \emph{Par não-adjacente}: após adicionar $\{w_1,z_1\}$ e $\{w_2,z_2\}$, as
    duas arestas restantes em $A \times Z$ são $\{w_1,z_2\}$ e $\{w_2,z_1\}$.
    \begin{itemize}
      \item Adicionar $\{w_1,z_2\}$: os seis vértices $\{a,b,c,w_1,z_1,z_2\}$ têm
        as arestas $w_1 z_1$ (adicionada), $w_1 z_2$ (sendo adicionada),
        $z_1 z_2$ (aresta $xy$ de $T_\ell$), e todas as arestas de
        $\Delta \cup \{w_1\} \cup \{z_1,z_2\}$ via as ligações ao triângulo.
        Verifica-se que esse conjunto forma $K_6$. $\bot$
      \item Analogamente para $\{w_2,z_1\}$. $\bot$
    \end{itemize}
\end{enumerate}
Portanto, \textbf{qualquer terceiro passo $U_\ell$ cria $K_6$}, o que é proibido
em Caso~1. Logo o número de passos $U_\ell$ neste caso é no máximo $\mathbf{2}$.

\medskip
\noindent\textit{Conclusão do Caso 1.}

Em qualquer construção de $K_5$-bloco $B$ sem $K_6$ e com $\maxdens(B) < 3$,
o primeiro passo $U_\ell$ é precedido por ao menos um passo $T_\ell$ (e,
possivelmente, por um passo $Q_\ell$ anterior), consumindo ao menos $1$ unidade
de orçamento antes do primeiro $U_\ell$. Além disso, a análise geométrica
da estrutura bipartida $A \times Z$ sobre o triângulo base $\Delta$ mostra
que qualquer terceiro passo $U_\ell$ cria obrigatoriamente um $K_6$. Portanto,
o número total de passos $U_\ell$ quando $B$ não contém $K_6$ é no máximo
$\mathbf{2}$, concluindo o Caso~1 e, juntamente com o Caso~2, a prova completa
da Afirmação~\ref{afirm:ulimit}. 
\medskip

\emph{[A completar.]}
\end{proof}




\vspace{1cm}

% ============================================================
\printbibliography
% ============================================================

\end{document}
